\section{Kernel tools}\label{sec:kernels}

We need two ways to retain mass in a shared environment. The stage
construction will require tuples of particles with separated endpoints
and a bound on rapid loss of their mass. The stationary-profile argument
will instead require a pair of particles whose signed gap does not
decrease, at a cost growing only polynomially with the traversed height.
We derive both conclusions from clipping and endpoint spread.

Given a probability \(\pi\) on \(k\)-tuples of sites in one layer, write
\[
\Gamma_h^\pi=\pi K_h^{\otimes k}
\]
using the same environment, as endpoint or prefix mass, and \({\rm sep}(\mathbf x)=\min_{i\ne j}|(x_i)_2-(x_j)_2|\). Write
\(\mathcal R_h(d')=\{\max_{s\le h,\ i,j}|Y_s^{(i)}-Y_s^{(j)}|\le d'\}\).
The event \(\mathcal R_h(d')\) concerns relative displacement differences
at first hits, not absolute positions; \(s\) is the integer height offset
from the indicated base. It reduces initial separation by at most \(d'\).
Such budgets add under concatenation at fresh layers. In each estimate
under \(P\), the initial law \(\pi\) is fixed independently of the
upper environment. The same estimate therefore applies conditionally
to a lower-layer-measurable initial law at a fresh layer.

We repeatedly use the following elementary composition rule. Retain a
restricted submeasure of \(\Gamma_h^\pi\), normalize its endpoint law,
and continue with a local kernel from that law. The product of the two
retained masses is then available in the longer original-floor kernel.
This follows from the quenched Markov property at ordinary first hits,
also for a product of walks: local no-drop conditions add restrictions
to original-floor success. Taking prefixes similarly gives lower bounds
at shorter horizons. Every finite-height mass, restriction, and
normalization used below depends only on rows below its top; it can be
computed by summing countably many finite path weights.

\begin{lemma}[Clipping from an initial distribution]\label{kernel-initial-clipping}
For fixed \(k,T,\eta,p'>0\), with integer \(k\ge2\), there is a
positive constant \(d\) such that, for every sufficiently large \(r\),
every \(H\le Tn(r)\), and every deterministic initial probability
\(\pi\),
\[
P\big(\Gamma_H^\pi(\mathcal R_H(\eta r))\ge d\big)>1-p'.
\]
The analogous single-walk clipping bound holds from any initial
probability, with the deterministic slope supplied by \eqref{eq:24}.
\end{lemma}

\begin{proof}
Apply \eqref{eq:24} with tolerance \(\eta/2\) and a sitewise failure
probability \(\epsilon\). For every tuple, all \(k\) clipped kernels
are good except with probability at most \(k\epsilon\). Hence the
expected \(\pi\)-fraction of bad tuples is at most \(k\epsilon\), and
with probability at least \(1-2k\epsilon\) at least half of the tuple
mass is good. On that event the product of clipped masses gives
\(d=d_0^k/2\). Choose \(2k\epsilon<p'\). The single-walk assertion is
the same argument with one coordinate.
\end{proof}

When a shorter-horizon test is bounded using clipping through a longer
horizon, the test itself still uses only the actual shorter kernel.
All these uniform bounds apply conditionally to a past-measurable
initial probability at a fresh layer.

\subsection{Creating separated seeds}

\begin{proposition}[Separated seeds]\label{kernel-seeds}
For each fixed \(k\ge2,\ p'>0\), there are constants \(C_*,c_0,g_0>0\) such that for all sufficiently large integer \(H\), putting \(r=r_{H/C_*}\),
\[
P\big(\Gamma_H^\pi\{{\rm sep}(\text{endpoints})\ge c_0 r\}\ge g_0\big)>1-p'
\tag{30}\label{eq:30}
\]
from arbitrary \(\pi\), without initial separation.
\end{proposition}

\begin{proof}
\smallskip\noindent\textbf{Splitting one starting mass.}
First consider a single prescribed start. We will produce \(k\) ordered
endpoint groups, with support gaps of order \(r\), each carrying a
fixed positive mass. Use \(J_0\) blocks of length
\(h=\lfloor n(r)\rfloor\), followed by a remainder, to reach
\(H=C_*n(r)\), where \(J_0\) is a large fixed integer and
\(C_*=2J_0+2\). Clipping through \(C_*n(r)\) supplies a common
deterministic slope \(s_r\) and, with arbitrarily high environment
probability, fixed positive mass within \(\eta r\) of that line.
The tolerance \(\eta\) will be small compared with \(1/J_0\).

Endpoint spread, Equation~\eqref{eq:13}, supplies a designated sign
of a deviation of size at least \(cr\) from \(hs_r\), with
single-walk \(P^+\)-probability at least \(c>0\). These constants
are independent of the clipping parameters. The corresponding raw
\(K_h\)-mass has expectation at least \(pc\), by using paths on
\(D\), and is bounded by one. It therefore exceeds a fixed positive
threshold on an environment event of fixed positive probability. The
same reasoning applies after averaging over any starting probability.

Maintain the groups as submeasures of the walked mass, normalizing
within each group for its next test. If there are fewer than \(k\)
groups and the designated sign is positive, take the rightmost group
and a second-coordinate median. Preserve its half at or below the
median by central clipping. From the half at or above the median,
retain the designated jumping mass whenever it exceeds its fixed
threshold; these endpoints form a new rightmost group. Preserve all
other groups by central clipping as well. For a negative designated
sign, perform the reflected operation at the leftmost group.

Both median halves have mass fraction at least one half. A median atom
may be used in both tests. In the positive-sign case, the retained
central endpoints have an upper bound given by the median plus \(hs_r+\eta r\), while the jumping
endpoints have a lower bound given by the median plus \(hs_r+cr\).
The negative-sign case has the reflected inequalities. Thus the new groups have disjoint supports with a gap of at least
\((c-\eta)r\). Every old gap loses at most \(2\eta r\) under
central propagation. Once there are \(k\) groups, only their central
preservation is needed.

The initial-distribution clipping lemma makes all the central
preservations succeed simultaneously with arbitrarily high conditional
probability. Abort if a preservation fails. Subtracting this small
error from the jump event shows that, whenever a split is still needed,
preservation together with a split has conditional probability at least
a fixed \(s_*>0\), independent of \(\eta,J_0\). The operations are
measurable: take full restrictions to the indicated events and choose,
for example, lower integer medians.

Choose \(J_0\) large, and then choose the preservation errors small
compared with \(1/J_0\). The probability of no abort but fewer than
\(k-1\) splits in these blocks is at most
\[
\sum_{i<k-1}\binom{J_0}{i}(1-s_*)^{J_0-i}.
\]
Indeed, on such a pattern a split is needed at every step, and each
prescribed nonsplit without abortion costs at most \(1-s_*\) by
conditioning on the previous blocks. Apply central clipping once more
over the remaining height, which is at most \(C_*n(r)\). Taking
\(\eta\) sufficiently small after fixing \(c,J_0\) leaves final
gaps at least \(c'r>0\). The product of the finitely many retained
mass fractions gives a fixed positive lower bound for each group in
the original single-walk kernel \(K_H\). Every test and continuation
uses only layers below its actual endpoint.

\smallskip\noindent\textbf{Choosing separated endpoints.}
For any starting tuple whose \(k\) starts each have these \(k\) options,
the product kernel gives fixed positive mass to separated endpoint
tuples. Choose \(2c_0<c'\) and expose endpoints successively. Every
earlier endpoint excludes at most one group at the next start, since
two distinct groups are separated by more than \(2c_0r\). Thus at the
\(i\)th choice at least one of the \(k\) groups remains available.
Each choice costs only its fixed positive group-mass lower bound.
Averaging the favorable fraction over \(\pi\), as in
Lemma~\ref{kernel-initial-clipping}, proves \eqref{eq:30}.
\end{proof}

\subsection{Rapid loss bound}

Fix \(\lambda\in(0,1]\) satisfying \eqref{eq:4} and write
\[
q_\pi=\int\prod_{i=1}^k q(x_i)\,\pi(d\mathbf x).
\]

\begin{lemma}[A separated inverse-moment bound]\label{kernel-inverse-moment}
For each fixed positive \(B_*<\infty\) and integer \(k\), if initial
separation on the support of \(\pi\) is sufficiently large, depending
on \(B_*,k\), then
\[
P(q_\pi<2s)\le e^{C k}s^\lambda,\qquad s\ge e^{-B_*},
\tag{31}\label{eq:31}
\]
with \(C\) independent of \(B_*,k\).
\end{lemma}

\begin{proof}
Truncate each \(q\) upward to \(q\vee e^{-B_*}\), changing the product
by at most \(e^{-B_*}\): if a factor changes, the entire new product is
at most this quantity. By bounded product-field mixing at mutually
separated sites, the expectation of the inverse \(\lambda\)-power of
the truncated product is bounded by \(2(Eq(0)^{-\lambda})^k\) for
sufficiently large separation. To see the required uniformity, at fixed
\(k,B_*\) approximate the bounded single-site functional in \(L^1(P)\)
by one using finitely many rows and having the same bound. Sufficient
separation makes all its translates disjoint. On the event on the left
of \eqref{eq:31}, the truncated product integral is \(<3s\). Jensen's
inequality for its negative power and Markov's inequality give
\eqref{eq:31}.
\end{proof}

\begin{proposition}[Rapid loss is unlikely]\label{kernel-rapid-loss}
Suppose \(\pi\) starts with \({\rm sep}\ge G\) and \(0<d'\le G/2\), with \(G-d'\) sufficiently large for \eqref{eq:31}. For any positive integer \(L'\) and \(kL'\log(1/\kappa)\le x\le B_*\), the loss \(Z=-\log\Gamma_H^\pi(\mathcal R_H(d'))\) satisfies
\[
P(Z>x)\le
\left(e^{Ck-\lambda x/L'} + C k e^{x/L'}\,H/n(d'/(2L'))\right)^{L'}.
\tag{32}\label{eq:32}
\]
Here \(C\) can be fixed independently of \(k,L',B_*,H,d'\), given the
stated separation condition.
\end{proposition}

\begin{proof}
From the current layer in a subdivision, monitor the local path mass
with budget \(d'/L'\), stopping at the first height where it is
\(<s:=\kappa^{-k}e^{-x/L'}\). The assumptions give
\(e^{-B_*}\le s\le1\). At any fresh start whose separation is still
at least \(G-d'\), the probability of such a threat is bounded by the
parenthesis in \eqref{eq:32}. Indeed the raw mass of all \(k\) infinite
\(D\)-successes is exactly the current \(q\)-product integral. Its
part violating the budget within \(H\) layers has expectation at most
\(CkH/n(d'/(2L'))\), by \eqref{eq:8}--\eqref{eq:9} and a union over
single walks on \(D\) with a median deviation greater than
\(d'/(2L')\). The other path factors are bounded by one and can be
dropped. If the product integral is at least \(2s\) and the violating
mass at most \(s\), enough mass satisfies the budget at every first
hit to prevent a threat. Lemma~\ref{kernel-inverse-moment} and Markov's
inequality now give the claimed bound; the factor
\(\kappa^{-k\lambda}\) is absorbed into \(e^{Ck}\). Neither \(D\) nor
\(q\) is used to choose the stopping layer: the actual test is the
finite-height mass, measurable from rows below that layer.

At a threat, extend the passing mass from one layer earlier by a direct
upward step for every particle; at offset zero the passing mass is one.
This preserves relative displacement errors and retains mass at least
\(\kappa^ks=e^{-x/L'}\). Normalize at the actual stopping layer and
restart with the fresh upper environment. A final piece with no threat
also has mass at least \(e^{-x/L'}\).

If the horizon finishes within \(L'\) pieces, concatenation gives
original mass at least \(e^{-x}\), with total relative-motion budget
at most \(d'\). Otherwise each of the first \(L'\) pieces has
detected a threat before completion. The conditional threat estimate
applies at every restart: only lower rows have been revealed, and
separation remains at least \(G-d'\). Multiplying these conditional
bounds gives Equation~\eqref{eq:32}.
\end{proof}

\subsection{An outward-order kernel}

For any same-layer sites \(x,y\) put
\[
p_H^\to(x,y)=(K_H(x)\otimes K_H(y))\{Y_H^{(2)}\ge Y_H^{(1)}\}.
\]
The inequality concerns movement displacements; the two quenched paths
use the same environment. Thus it asks that their signed physical gap
not decrease from its initial value.

\begin{proposition}[Outward order at polynomial cost]\label{kernel-outward-order}
There are fixed \(A_{\to},c>0\) and a nonnegative random function
\(\mathcal Z(x,y)\) of the upper environment such that, simultaneously
in integer \(H\ge1\),
\[
-\log p_H^\to(x,y)\le A_{\to}\log(1+H)+\mathcal Z(x,y),\qquad
\sup_{x,y} E\exp(c\mathcal Z(x,y))<\infty.
\tag{33}\label{eq:33}
\]
Every cap \(\mathcal Z\wedge B'\), \(0<B'<\infty\), is approximable in \(P\)-probability uniformly in \(x,y\) by bounded functions of fixed-radius upper row neighborhoods of these sites.
\end{proposition}

\begin{proof}
We construct one sequence of retained endpoint laws and use it at every
target height. A stage starts with an extra signed gap, or buffer,
between the walks. It usually increases this buffer by a fixed factor;
on failure it reduces the buffer by a smaller factor. The resulting
positive logarithmic growth limits the number of stages through height
\(H\) to order \(\log(1+H)\), apart from an exponentially integrable
random error. Each stage costs only a fixed amount of logarithmic mass.

\paragraph{Endpoint tests and retention.}
Keep the original signed gap \(z_0=y_2-x_2\) fixed, even when it is
negative. First prescribe \(\lceil R_0\rceil\) positive lateral
steps for walk 2 and one upward step for each walk. This bounded
uniformly elliptic script supplies a buffer \(R_0\), to be fixed
large enough below, at a fixed positive mass cost.

At a stage base, the current pair probability \(\pi\) is supported
on signed gaps at least \(z_0+R\), where \(R\ge R_0\). Plan
\(h=\lfloor n(R)\rfloor\) layers. For constants
\(a>0\), \(0<\ell'<1\), and \(g_*>0\) to be chosen, test
\[
\Gamma_j^\pi\{\text{signed endpoint gap}\ge z_0+(1-\ell')R\}
   \ge g_*
\qquad(1\le j\le h).
\]
Stop at the first failing test. If all pass, also test at \(h\)
whether mass at least \(g_*\) has gap at least \(z_0+(1+a)R\).
These are tests of the full kernels from the same stage base. We do not
intersect their endpoint events as path restrictions: a completed stage
uses its retained final law, while a target height inside that stage
uses only the test at that one height.

On success retain and normalize the growing endpoint mass, and replace
\(R\) by \((1+a)R\). On failure replace \(R\) by
\((1-\ell')R\). At a first traversal failure, take the passing
mass from the preceding layer and extend each walk by one direct upward
step. This preserves its signed gap and supplies mass at least
\(\kappa^2g_*\) at the stopping layer; at offset zero the passing
mass is one. If only the final growth test fails, its passing traversal
test already supplies this much mass. Normalize the retained law.
Whenever the new buffer is below \(R_0\), restore it by bounded
positive lateral steps for walk 2 followed by an upward step for both
walks. The buffer before restoration is at least
\((1-\ell')R_0\), so the script has uniformly bounded length.

Every test, retained law, and restoration uses only departure rows
below its terminal layer. Conditional on this information, the next
stage therefore starts in fresh iid layers.

\paragraph{A uniform chance of buffer growth.}
We choose \(a\) and a success probability lower bound \(s_*>0\)
independently of small \(\ell'\). Clip both walks at scale \(R\)
about the common deterministic slope in Equation~\eqref{eq:24}, with
tolerance \(\delta R\), where \(2\delta<\ell'\). The endpoint
spread estimate~\eqref{eq:13}, applied to the same deterministic
center \(hs_R\), gives one sign of a deviation of size at least
\(c_1R\) with single-walk annealed \(P^+\)-probability at least a
fixed positive constant. The constants do not depend on the clipping
parameters. The favorable sign may depend on \(R\): use a positive
deviation for walk 2, or a negative deviation for walk 1. Either choice
increases their signed gap when the other walk stays near the center.

The raw jump mass has a fixed positive expectation, by restricting to
\(D\)-paths, and is bounded by one. Thus it is at least some
\(j_0>0\) on an environment event of probability at least
\(c_0>0\), uniformly over starts and large \(R\). Choose the
sitewise clipping errors with sum less than \(c_0/2\), and let
\(d>0\) be their common retained mass bound. For each deterministic
starting pair \(\mathbf x\), the event \(A_{\mathbf x}\) that
this jump mass is at least \(j_0\) and both clipped masses are at
least \(d\) then satisfies
\[
P(A_{\mathbf x})\ge c_0/2.
\]
This uses subtraction of the two clipping errors, with no independence
between the events.

For an arbitrary current pair law, set
\(F=\int\mathbf1_{A_{\mathbf x}}\,\pi(d\mathbf x)\).
Since \(EF\ge c_0/2\) and \(0\le F\le1\),
\[
P(F\ge c_0/4)\ge c_0/4.
\]
On this event, central clipping for the favorable pairs supplies at
least \(Fd^2\) to every traversal test. At the terminal height,
the favorable jump and the central restriction for the other walk
supply at least \(Fj_0d\) to the growth test. The terminal paths
need not be the paths used to prove the traversal tests; the tests
were defined separately for this reason. Choose \(\delta<c_1/2\),
\(a<c_1/2\), and
\[
g_*\le(c_0/4)\min(d^2,j_0d).
\]
Then every stage succeeds with conditional probability at least a
fixed \(s_*>0\), independently of small \(\ell'\). The constants
\(g_*\) and \(R_0\) may depend on the chosen \(\ell'\).
All choices can be made deterministic at the countable set of buffer
values reached by the procedure.

\paragraph{The number of stages through a given height.}
Fix \(\ell'\) small enough that a proportion \(s_*/2\) of
successes gives positive logarithmic growth. Let \(S_m\) count
successes in the first \(m\) main stages, and put
\[
W=\sup_{m\ge1}(s_*m/2-S_m)_+.
\]
The conditional success lower bound implies a uniform exponential tail
for \(W\). Indeed, for sufficiently small \(\theta>0\),
\[
\exp\{\theta(s_*m/2-S_m)\}
\]
is a nonnegative supermartingale: the conditional multiplier is at most
\(e^{\theta s_*/2}(1-s_*+s_*e^{-\theta})\le1\).
Its maximal inequality gives \(P(W>u)\le e^{-\theta u}\).

Success multiplies the buffer by \(1+a\), failure by \(1-\ell'\),
and restoration only increases it. Consequently, after \(m\)
completed main stages and their restorations,
\[
\log R\ge c_2m-C_2(W+1)
\]
for fixed positive constants \(c_2,C_2\). If \(H'\) is the
height already consumed, we also have \(R\le C(H'+1)\).
A successful increase from \(R_{\rm old}\) uses its full planned
length \(\lfloor n(R_{\rm old})\rfloor\ge cR_{\rm old}\), by
Equation~\eqref{eq:7} and a sufficiently large \(R_0\). Failures
reduce the buffer, and restorations reset it to \(R_0\).
Thus at any integer target height \(H\), at most
\[
C\bigl(\log(1+H)+W+1\bigr)
\]
main stages have been completed, and at most one more is unfinished.

Concatenate the retained masses through \(H\). If the last stage is
unfinished, use its traversal test or its failure retention at the
actual offset. If the target is the endpoint of a one-layer restoration,
its bounded script supplies the continuation instead. The endpoint gap
stays above the original baseline, so
these paths contribute to \(p_H^\to(x,y)\). Every completed stage,
including a possible restoration, and the last partial stage cost a
fixed logarithmic amount. Hence, simultaneously for all \(H\ge1\),
\[
-\log p_H^\to(x,y)
\le C_3\log(1+H)+C_3(W+1)
\]
with a constant uniform in the starting pair.

\paragraph{An exponentially integrable local remainder.}
Take \(A_\to=2C_3\), and define
\[
\mathcal Z(x,y)=\sup_{H\ge1}
\bigl[-\log p_H^\to(x,y)-A_\to\log(1+H)\bigr]_+.
\]
The preceding bound gives \(\mathcal Z\le C_3(W+1)\), proving
the exponential moment in Equation~\eqref{eq:33}. It also gives
\[
-\log p_H^\to(x,y)-A_\to\log(1+H)
\le C_3(W+1)-C_3\log(1+H).
\]
On \(\{W\le M\}\), only a deterministic bounded range of heights
can therefore contribute to the supremum. The probability of its
complement tends to zero uniformly in \(x,y\).

For a bounded height range, kernels can be uniformly approximated by
bounded-length paths. Indeed, from any site inside a finite-height
strip, a fixed number of consecutive upward steps forces exit and has
uniformly positive probability. The strip exit time therefore has a
geometric block bound, uniformly in the environment and horizontal
start. Restricting both paths to a bounded number of steps changes
their pair kernel mass by a uniformly small amount and uses only
fixed-radius upper-row neighborhoods of the two starts. Those
neighborhoods may overlap; the truncation estimate is unchanged.
Finally,
\[
p_H^\to(x,y)\ge\kappa^{2H},
\]
by the two straight upward paths. Logarithms are therefore uniformly
continuous on this bounded range of heights. Take the finite maximum,
then cap at \(B'\), and let the two truncation errors tend to zero.
This proves the uniform local approximation of
\(\mathcal Z\wedge B'\).
\end{proof}
